% =====================================================================
%  INFORME INSTITUCIONAL — CENTRO DE INTERÉS "ConectaTE"
%  Solicitud de presupuesto: $2.000.000 COP
%  Impresión de las guías de la plataforma ConectaTE
%  + Libro de caricaturas "Salomé y los exoplanetas"
%  + Revista Cosmotec (3 fascículos · semillero institucional)
%  ---------------------------------------------------------------------
%  Institución Educativa Sor María Juliana
%  Cartago, Valle del Cauca · Colombia
%  Docente responsable: PhD. Álvaro Cárdenas Orozco
%  Compilar con: xelatex informe_conectate_centro_interes.tex
% =====================================================================

\documentclass[12pt, letterpaper]{article}

% --- Idioma y codificación ---
\usepackage[spanish, es-tabla, es-nodecimaldot]{babel}

% --- Geometría / márgenes ---
\usepackage[top=2.5cm, bottom=2.5cm, left=2.8cm, right=2.5cm, headheight=24pt]{geometry}

% --- Colores institucionales SMJ ---
\usepackage{xcolor}
\definecolor{smjAzul}     {HTML}{0D1B3E}
\definecolor{smjVino}     {HTML}{6B1A2A}
\definecolor{smjDorado}   {HTML}{C9A84C}
\definecolor{smjGrisClaro}{HTML}{F4F4F4}
\definecolor{smjGrisLinea}{HTML}{D0D4DC}

% --- Fuentes (XeLaTeX) ---
\usepackage{fontspec}
\setmainfont{Times New Roman}

% --- Gráficos ---
\usepackage{graphicx, float}
\usepackage{tikz}
\usetikzlibrary{shapes, arrows.meta, positioning, calc, mindmap, trees}

% --- Tablas ---
\usepackage{booktabs, array, tabularx, longtable, multirow}
\usepackage{colortbl}
\usepackage{ragged2e}
\newcolumntype{L}[1]{>{\RaggedRight\arraybackslash}p{#1}}
\newcolumntype{C}[1]{>{\Centering\arraybackslash}p{#1}}
\newcolumntype{R}[1]{>{\RaggedLeft\arraybackslash}p{#1}}

% --- Listas, formato general ---
\usepackage{enumitem}
\setlist{itemsep=2pt, topsep=4pt}
\usepackage{parskip}
\usepackage{microtype}
\usepackage{setspace}
\onehalfspacing

% --- Encabezados de sección ---
\usepackage{titlesec}
\titleformat{\section}
  {\Large\bfseries\color{smjAzul}}
  {\thesection.}{0.6em}{}
  [{\color{smjDorado}\titlerule[1pt]}]
\titleformat{\subsection}
  {\large\bfseries\color{smjVino}}
  {\thesubsection}{0.6em}{}

% --- Encabezado y pie ---
\usepackage{fancyhdr}
\pagestyle{fancy}
\fancyhf{}
\fancyhead[L]{\footnotesize\color{smjAzul}\textit{Centro de Interés ConectaTE}}
\fancyhead[C]{\footnotesize\color{smjVino}\textit{I.E. Sor María Juliana}}
\fancyhead[R]{\footnotesize\color{smjAzul}Pág. \thepage}
\fancyfoot[C]{\footnotesize\color{gray}Ecosistema STEAM+ · PhD. Álvaro Cárdenas Orozco}
\renewcommand{\headrulewidth}{0.4pt}
\renewcommand{\headrule}{\hbox to\headwidth{\color{smjDorado}\leaders\hrule height \headrulewidth\hfill}}
\renewcommand{\footrulewidth}{0.2pt}

% --- Cajas ---
\usepackage[most]{tcolorbox}

\newtcolorbox{cajaResumen}{
  enhanced,
  colback=smjAzul!4, colframe=smjAzul,
  boxrule=0.6pt, arc=2mm,
  left=10pt, right=10pt, top=8pt, bottom=8pt,
  fonttitle=\bfseries\color{white},
  coltitle=white, colbacktitle=smjAzul,
  title=Resumen ejecutivo
}

\newtcolorbox{cajaCita}{
  enhanced, breakable,
  colback=smjGrisClaro, colframe=smjDorado!80!black,
  boxrule=0.4pt, arc=1mm,
  left=12pt, right=12pt, top=6pt, bottom=6pt
}

\newtcolorbox{cajaTotal}{
  enhanced,
  colback=smjVino!8, colframe=smjVino,
  boxrule=0.8pt, arc=2mm,
  left=10pt, right=10pt, top=8pt, bottom=8pt
}

\newtcolorbox{cajaConcepto}[1]{
  enhanced, breakable,
  colback=smjAzul!5, colframe=smjAzul,
  boxrule=0.5pt, arc=2mm,
  left=10pt, right=10pt, top=6pt, bottom=6pt,
  fonttitle=\bfseries\color{white},
  coltitle=white, colbacktitle=smjAzul,
  title=#1
}

% --- Hyperref al final ---
\usepackage[hidelinks]{hyperref}
\hypersetup{
  pdftitle={Informe Centro de Interés ConectaTE — Solicitud de Presupuesto},
  pdfauthor={PhD. Álvaro Cárdenas Orozco},
  pdfsubject={ConectaTE — Ecosistema STEAM+ · I.E. Sor María Juliana},
  pdfkeywords={ConectaTE, STEAM+, MILC, CIR Sormatronik, Betelgeuse, exoplanetas, Salomé, Cartago}
}

% --- Comandos de utilidad ---
\newcommand{\pesos}[1]{\$\,#1}
\newcommand{\rubrica}[1]{\textcolor{smjVino}{\textbf{#1}}}
\newcommand{\conectate}{\textbf{ConectaTE}}

% =====================================================================
\begin{document}

% ============================== PORTADA ==============================
\thispagestyle{empty}
\begin{center}
{\color{smjAzul}\rule{\linewidth}{2.2pt}}\\[0.2cm]
{\large\textsc{\color{smjAzul} República de Colombia}}\\[0.05cm]
{\large\textsc{\color{smjAzul} Institución Educativa Sor María Juliana}}\\[0.05cm]
{\normalsize\color{smjVino} Cartago — Valle del Cauca}\\[0.10cm]
{\small\color{smjDorado!80!black} Centro de Interés Institucional · Ecosistema STEAM+}\\[0.25cm]
{\color{smjDorado}\rule{\linewidth}{1pt}}\\[1.3cm]

% Tipo de documento
{\large\bfseries\color{smjVino} INFORME INSTITUCIONAL}\\[0.2cm]
{\small\color{smjAzul} Presentación del Centro de Interés y solicitud de presupuesto}\\[1.0cm]

% Título principal
{\Huge\bfseries\color{smjAzul} ConectaTE}\\[0.25cm]
{\Large\bfseries\color{smjAzul} Centro de Interés Institucional}\\[0.10cm]
{\large\bfseries\color{smjVino} Ciencia, Tecnología, Educación e Investigación Escolar}\\[0.10cm]
{\normalsize\itshape\color{smjVino}Un ecosistema pedagógico STEAM+ de la I.E.\ Sor María Juliana}\\[0.9cm]

% Subtítulo
\begin{minipage}{0.90\textwidth}\centering
{\itshape\color{smjAzul}\normalsize
Solicitud de financiación por valor de \pesos{2.000.000} (dos millones de pesos m/cte) para la \mbox{impresión}, encuadernación y distribución de tres productos editoriales del Centro de Interés ConectaTE: las \textbf{guías didácticas} de su plataforma educativa institucional, el \textbf{libro de caricaturas científicas «Salomé y los exoplanetas»} y la \textbf{revista Cosmotec} (dos fascículos en reedición y uno inédito), cuya circulación está orientada a sostener económicamente al semillero institucional.}
\end{minipage}\\[1.1cm]

{\color{smjDorado}\rule{0.5\linewidth}{0.5pt}}\\[0.6cm]

% Datos generales
\begin{tabular}{r l}
\textbf{\color{smjAzul} Docente responsable:} & PhD. Álvaro Cárdenas Orozco \\[2pt]
\textbf{\color{smjAzul} Centro de Interés:}   & ConectaTE — Ecosistema STEAM+ institucional \\[2pt]
\textbf{\color{smjAzul} Semilleros articulados:} & CIR Sormatronik (robótica) · Betelgeuse (astronomía) \\[2pt]
\textbf{\color{smjAzul} Grados beneficiados:} & 6.°{}\,, 7.°{}\,, 8.°{}\,, 9.°{}\,, 10.°{}{} y 11.°{} \\[2pt]
\textbf{\color{smjAzul} Dirigido a:}          & Rectoría — Consejo Directivo \\[2pt]
\textbf{\color{smjAzul} Vigencia:}            & Año lectivo 2026 \\
\end{tabular}\\[1.1cm]

{\normalsize\color{smjVino} Cartago, \today}\\[0.3cm]
{\color{smjAzul}\rule{\linewidth}{2.2pt}}
\end{center}

\newpage

% ============================ RESUMEN ============================
\section*{Resumen ejecutivo}
\addcontentsline{toc}{section}{Resumen ejecutivo}

\begin{cajaResumen}
\conectate{} es el \textbf{Centro de Interés Institucional} de la I.E. Sor María Juliana orientado a integrar ciencia, tecnología, educación e investigación escolar mediante experiencias \textbf{STEAM+} contextualizadas. No es un semillero extracurricular, sino un \textit{ecosistema pedagógico} que articula robótica, astronomía, pensamiento computacional, inteligencia artificial, creatividad, investigación y formación ciudadana, y que recoge la trayectoria de los semilleros institucionales \textbf{CIR Sormatronik} y \textbf{Betelgeuse}.

De ConectaTE como centro madre se derivan productos pedagógicos, científicos y tecnológicos. \textbf{Tres} de ellos son objeto de esta solicitud: (i) las \textbf{180 guías didácticas} y los \textbf{proyectos integradores} de la plataforma educativa institucional, (ii) el \textbf{libro de caricaturas científicas «Salomé y los exoplanetas»} derivado del semillero Betelgeuse, y (iii) la \textbf{revista institucional Cosmotec} —dos fascículos en reedición y uno inédito— cuya circulación bajo esquema de aporte simbólico está destinada a \textit{financiar la operación del semillero institucional} (insumos de prototipos, transporte a ferias, materiales de laboratorio). Esta solicitud por \textbf{\pesos{2.000.000} COP} cubre la impresión, encuadernación y distribución de los tres productos, y activa un mecanismo de autosostenibilidad económica del semillero. Beneficia directamente a \textbf{180 estudiantes} de 6.°{} a 11.°{} en el año lectivo 2026.
\end{cajaResumen}

% ============================ TABLA DATOS ============================
\section{Datos generales del proyecto}

\renewcommand{\arraystretch}{1.30}
\begin{tabularx}{\linewidth}{>{\bfseries\color{smjAzul}}p{4.6cm} X}
\toprule
Nombre del proyecto      & Centro de Interés Institucional \textbf{ConectaTE} — Ecosistema STEAM+ \\
Naturaleza               & Centro de interés institucional de innovación educativa, investigación escolar y producción editorial científica \\
Institución              & Institución Educativa Sor María Juliana — Cartago, Valle del Cauca \\
Docente responsable      & PhD. Álvaro Cárdenas Orozco \\
Líneas articuladas       & Robótica · Astronomía · Pensamiento computacional · IA · Tecnología · Creatividad · Investigación · Formación ciudadana \\
Semilleros institucionales & CIR Sormatronik (robótica) · Betelgeuse (astronomía) \\
Población beneficiaria   & Estudiantes de 6.°{}, 7.°{}, 8.°{}, 9.°{}, 10.°{} y 11.°{} (aprox.\ 180 estudiantes) \\
Vigencia                 & Año lectivo 2026 (4 periodos académicos) \\
Productos objeto de la solicitud & (A) 180 guías didácticas de la plataforma ConectaTE \par (B) Libro «Salomé y los exoplanetas» — semillero Betelgeuse \par (C) Revista \textbf{Cosmotec} — 3 fascículos (2 reediciones + 1 inédito) para financiamiento del semillero \\
Enfoque                  & STEAM+ contextualizada · Ciencia ciudadana · Cultura institucional \\
Valor solicitado         & \rubrica{\pesos{2.000.000} COP — dos millones de pesos m/cte} \\
Fuente sugerida          & Fondo de Servicios Educativos / Rectoría / Aliado externo \\
\bottomrule
\end{tabularx}

% ============================ PRESENTACIÓN ============================
\section{ConectaTE como Centro de Interés}

\conectate{} puede definirse como un \textbf{Centro de Interés Institucional} orientado a integrar ciencia, tecnología, educación e investigación escolar mediante experiencias STEAM+ contextualizadas, donde las y los estudiantes aprenden creando, investigando, resolviendo problemas reales y participando en escenarios de innovación, ciencia ciudadana y transformación comunitaria.

En este sentido, \conectate{} no es solamente un semillero extracurricular, sino un \textbf{ecosistema pedagógico} que articula robótica, astronomía, pensamiento computacional, inteligencia artificial, tecnología, creatividad, investigación y formación ciudadana. Su propósito es convertir el interés de las y los estudiantes por la ciencia y la tecnología en una \textit{ruta formativa permanente}, capaz de fortalecer competencias investigativas, pensamiento crítico, liderazgo, trabajo colaborativo, comunicación científica y solución de problemas del entorno.

Como Centro de Interés, \conectate{} funciona como un \textbf{laboratorio vivo de aprendizaje}, donde las estudiantes y los estudiantes de la Institución Educativa Sor María Juliana exploran preguntas, diseñan proyectos, construyen prototipos, participan en convocatorias, socializan resultados y conectan el currículo escolar con desafíos reales de la sociedad. Allí, el conocimiento deja de ser un contenido aislado y se convierte en una \textit{experiencia práctica, situada y significativa}.

\subsection{El enfoque STEAM+}

Su enfoque se fundamenta en la educación \textbf{STEAM+}, entendida como la integración de ciencia, tecnología, ingeniería, arte y matemáticas, ampliada con áreas complementarias indispensables para una formación integral: \textbf{ciudadanía, ambiente, comunicación, ética, emprendimiento e inclusión}. Por eso \conectate{} permite que los aprendizajes de Tecnología e Informática dialoguen con las ciencias naturales, la astronomía, la robótica, la programación, la inteligencia artificial, la lectura crítica, la escritura científica y la comprensión del territorio.

\begin{cajaConcepto}{STEAM+ en ConectaTE}
\textbf{S} Ciencia \quad \textbf{T} Tecnología \quad \textbf{E} Ingeniería \quad \textbf{A} Arte \quad \textbf{M} Matemáticas\\
\textbf{+} Ciudadanía · Ambiente · Comunicación · Ética · Emprendimiento · Inclusión
\end{cajaConcepto}

\subsection{Estrategia de gestión del conocimiento institucional}

\conectate{} se proyecta además como una \textbf{estrategia de gestión del conocimiento institucional}, porque:

\begin{itemize}[leftmargin=*]
  \item Recoge la trayectoria de los semilleros \textbf{CIR Sormatronik} (robótica e ingeniería) y \textbf{Betelgeuse} (astronomía y ciencia ciudadana).
  \item Consolida liderazgos estudiantiles y promueve la formación entre pares.
  \item Genera \textbf{productos educativos, científicos y tecnológicos} (plataformas, guías, libros, prototipos, dashboards, modelos pedagógicos) que pueden permear el currículo institucional.
  \item Transforma la curiosidad en investigación, la investigación en proyecto, el proyecto en experiencia formativa y la experiencia formativa en \textit{cultura institucional}.
\end{itemize}

En síntesis, \conectate{} como Centro de Interés es una \textbf{propuesta pedagógica de innovación educativa} que conecta tecnología y educación para formar estudiantes investigadoras, investigadores, creativas, creativos, críticas, críticos y comprometidos con su comunidad, capaces de imaginar soluciones y construir futuro desde la escuela pública.

% ============================ ECOSISTEMA ============================
\section{Ecosistema de productos derivados de ConectaTE}

Una característica esencial del Centro de Interés \conectate{} es que \textbf{la plataforma educativa y el modelo pedagógico son productos de ConectaTE, no a la inversa}. Es decir, primero existe el Centro de Interés como ecosistema vivo y, a partir de su trabajo sostenido, surgen los productos que materializan, comunican y replican lo que allí se aprende.

\renewcommand{\arraystretch}{1.30}
\begin{table}[H]
\centering
\small
\begin{tabularx}{\linewidth}{L{3.6cm} L{3.0cm} X}
\toprule
\rowcolor{smjAzul!12}
\textbf{Producto} & \textbf{Origen en ConectaTE} & \textbf{Descripción} \\
\midrule
\textbf{Plataforma educativa ConectaTE} & Línea de tecnología, IA y pensamiento computacional & Sitio web institucional \textit{offline-first} (Astro + TypeScript + Service Worker) que aloja las guías, proyectos integradores, laboratorios interactivos y bitácoras de los semilleros. \\
\textbf{180 guías didácticas} & Práctica de aula y producción editorial & Conjunto editorial de 30 guías por grado (6.°{} a 11.°{}), validadas por linter automático, conformes al contrato editorial institucional. \\
\textbf{15 proyectos integradores} & Trabajo de cierre de cada periodo & Proyectos transversales que articulan dos o más áreas en torno a un producto tecnológico verificable. \\
\textbf{Modelo pedagógico MILC} & Sistematización de la práctica & Modelo Institucional de Lectura Crítica desarrollado en el Centro de Interés, hoy aplicado a todas las guías de la plataforma. \\
\textbf{Libro «Salomé y los exoplanetas»} & Semillero Betelgeuse & Libro de caricaturas científicas que divulga la astronomía contemporánea desde la mirada de una niña afrocolombiana del Valle del Cauca. \\
\textbf{Revista Cosmotec} & Semillero institucional del Centro de Interés & Publicación periódica de divulgación científica y tecnológica producida por las y los estudiantes del semillero; su circulación está orientada a generar ingresos propios para sostener la operación del semillero. \\
\textbf{Prototipos del semillero CIR Sormatronik} & Semillero CIR Sormatronik & Prototipos de robótica e ingeniería desarrollados con y por las y los estudiantes. \\
\bottomrule
\end{tabularx}
\caption{Ecosistema de productos derivados del Centro de Interés ConectaTE.}
\end{table}

De este ecosistema, la presente solicitud focaliza \textbf{tres productos editoriales} con vocación de impacto inmediato y de sostenibilidad económica del Centro de Interés: (A) las \textit{guías didácticas} de la plataforma, (B) el \textit{libro de caricaturas} «Salomé y los exoplanetas» y (C) la \textit{revista Cosmotec} en sus tres fascículos (dos reediciones y uno inédito).

% ============================ JUSTIFICACIÓN ============================
\section{Justificación de la solicitud}

\subsection{Pertinencia pedagógica}

Llevar a formato físico los productos editoriales de \conectate{} responde a necesidades concretas observadas en aula:

\begin{itemize}[leftmargin=*]
  \item \textbf{Lectura profunda sin pantalla.} La evidencia neurocognitiva (Mangen, Wolf, Carr) muestra que la lectura sostenida en papel mejora la comprensión profunda frente a la lectura en pantalla. Las secciones de \textit{saber ancestral} y \textit{cierre filosófico} —ejes del modelo MILC— ganan profundidad en formato impreso.
  \item \textbf{Equidad de acceso.} Aunque la plataforma es \textit{offline-first}, no todas las estudiantes y todos los estudiantes tienen un dispositivo personal estable. El material impreso garantiza que nadie quede atrás por conectividad o equipo.
  \item \textbf{Cultura del cuaderno.} El modelo MILC exige escritura sistemática en el cuaderno físico; la guía impresa sobre el escritorio refuerza el hábito de \textit{leer–subrayar–escribir}.
  \item \textbf{Patrimonio institucional físico.} Cada cartilla y cada libro impreso queda como activo de la biblioteca y del archivo de la Sor María Juliana, disponible para futuras generaciones de estudiantes.
\end{itemize}

\subsection{Pertinencia científica y STEAM+}

El libro \textbf{«Salomé y los exoplanetas»} —producto editorial del semillero Betelgeuse— acerca a la comunidad estudiantil a una de las fronteras más vivas de la ciencia contemporánea: la detección, caracterización y habitabilidad de planetas fuera del Sistema Solar (más de 5\,500 exoplanetas confirmados; misiones Kepler, TESS y JWST). El recurso integra cuatro lenguajes en una sola pieza editorial:

\begin{itemize}[leftmargin=*]
  \item \textbf{Narrativa gráfica:} Salomé, una niña afrocolombiana del Valle del Cauca, descubre el método del tránsito observando el cielo desde la finca de su abuela.
  \item \textbf{Divulgación científica rigurosa:} curvas de luz, zona de habitabilidad, espectroscopía atmosférica, sistemas binarios.
  \item \textbf{Pensamiento computacional:} Salomé aprende a modelar tránsitos con código sencillo (Python y hoja de cálculo), articulado con las guías de datos e IA de la plataforma.
  \item \textbf{Diálogo ético-ciudadano:} el último capítulo pregunta qué significa habitar la Tierra si la vida puede existir en otros mundos.
\end{itemize}

Este producto es, además, una \textbf{declaración política y pedagógica}: que una niña afrocolombiana de Cartago piense la frontera del universo no es un detalle estético, es la afirmación cotidiana de que la ciencia se hace —también— desde la escuela pública del Valle del Cauca.

\subsection{Pertinencia institucional}

La solicitud se alinea con (i) el Proyecto Educativo Institucional de la Sor María Juliana, (ii) los Estándares Básicos de Competencias en Tecnología e Informática del MEN (Guía 30), (iii) los Derechos Básicos de Aprendizaje (DBA) en ciencias y tecnología, (iv) los lineamientos de \textit{ciencia ciudadana} y \textit{cultura digital} de la política pública nacional, y (v) la tradición institucional de los semilleros CIR Sormatronik y Betelgeuse. Formalizar a \conectate{} como Centro de Interés es reconocer institucionalmente lo que ya ocurre de hecho.

% ============================ OBJETIVOS ============================
\section{Objetivos}

\subsection{Objetivo general}

Fortalecer el Centro de Interés Institucional \conectate{} mediante la \textbf{producción, impresión y distribución} de dos de sus productos editoriales —las 180 guías didácticas de su plataforma y el libro de caricaturas «Salomé y los exoplanetas» del semillero Betelgeuse—, para consolidar el ecosistema STEAM+ de la I.E.\ Sor María Juliana y beneficiar de manera directa a las y los estudiantes de 6.°{} a 11.°{} durante el año lectivo 2026.

\subsection{Objetivos específicos}

\begin{enumerate}[label=\textbf{\color{smjAzul}OE\arabic*.}, leftmargin=*]
  \item \textbf{Imprimir} los kits maestros de las 180 guías didácticas de la plataforma ConectaTE, organizados por grado y periodo, para uso permanente en las aulas.
  \item \textbf{Producir} bitácoras de trabajo impresas para las y los estudiantes, con la selección de guías clave del año por grado.
  \item \textbf{Editar e imprimir} 60 ejemplares del libro «Salomé y los exoplanetas», dotando la biblioteca, las aulas y el archivo institucional.
  \item \textbf{Imprimir} 240 ejemplares de la revista \textbf{Cosmotec} —dos fascículos en reedición y uno inédito— y activar un esquema de aporte simbólico para que los ingresos sostengan al semillero del Centro de Interés.
  \item \textbf{Articular} los recursos físicos con la plataforma digital mediante códigos QR de complemento hacia laboratorios interactivos, dashboards y registros de los semilleros.
  \item \textbf{Visibilizar} el Centro de Interés ante la comunidad educativa mediante un acto de lanzamiento y un informe ejecutivo al Consejo Directivo.
\end{enumerate}

% ============================ POBLACIÓN ============================
\section{Población beneficiada}

\renewcommand{\arraystretch}{1.20}
\begin{table}[H]
\centering
\small
\begin{tabularx}{\linewidth}{C{1.4cm} L{2.0cm} C{1.8cm} C{1.6cm} X}
\toprule
\rowcolor{smjAzul!12}
\textbf{Grado} & \textbf{Cobertura} & \textbf{Estudiantes (aprox.)} & \textbf{Guías producidas} & \textbf{Énfasis STEAM+ del Centro de Interés} \\
\midrule
6.°{}  & Directa & 30 & 30 & Alfabetización digital, cultura del cuaderno, primeros pasos en pensamiento computacional \\
7.°{}  & Directa & 30 & 30 & Algoritmos, oficina avanzada, robótica introductoria (línea CIR Sormatronik) \\
8.°{}  & Directa & 30 & 30 & Datos, visualización, electrónica básica, IoT \\
9.°{}  & Directa & 30 & 30 & Programación con Python, web, modelado científico, semillero Betelgeuse \\
10.°{} & Directa & 30 & 30 & Inteligencia artificial, ética digital, robótica avanzada \\
11.°{} & Directa & 30 & 30 & Proyecto integrador final, investigación escolar, comunicación científica \\
\midrule
\multicolumn{2}{r}{\textbf{Total estudiantes directos}} & \textbf{180} & \textbf{180} & \\
\bottomrule
\end{tabularx}
\caption{Cobertura del Centro de Interés ConectaTE por grado.}
\end{table}

\textbf{Beneficiarios indirectos:} aproximadamente 360 acudientes, 28 docentes de distintas áreas (las guías son recurso transversal), la comunidad lectora de la biblioteca institucional y las redes pedagógicas externas (red de docentes investigadores del Valle del Cauca).

% ============================ COMPONENTES ============================
\section{Productos editoriales objeto de la solicitud}

\subsection{Producto A · Guías didácticas de la plataforma ConectaTE}

\textbf{Insumo digital existente:} 180 guías didácticas (30 por grado, de 6.°{} a 11.°{}) estructuradas en archivos \texttt{YAML} como única fuente de verdad, compiladas a PDF mediante pipeline propio (\texttt{make guia-build}), validadas por linter automático (\texttt{make guia-lint}) y conformes al modelo editorial MILC.

\textbf{Productos físicos a financiar:}
\begin{itemize}[leftmargin=*]
  \item \textbf{24 kits maestros de aula} (cartillas anilladas, una por grado-periodo, full color, portada plastificada): 6 grados $\times$ 4 periodos.
  \item \textbf{Bitácoras estudiantiles} con selección de guías clave por grado (b/n, perforadas, archivables): 180 ejemplares para el universo estudiantil.
  \item Material complementario: separadores, fundas, rotulado institucional.
\end{itemize}

\subsection{Producto B · Libro «Salomé y los exoplanetas» (semillero Betelgeuse)}

\textbf{Insumo digital existente:} guion y storyboard del libro elaborados por el docente responsable junto con el semillero Betelgeuse, articulados con las guías de astronomía, datos e IA del Producto A.

\textbf{Productos físicos a financiar:} 60 ejemplares impresos, formato A5 (15{,}0\,\(\times\)\,21{,}0\,cm), entre 56 y 64 páginas, full color, encuadernación rústica cosida con solapas, papel propalcote 115\,g (interior) y 250\,g (portada). Destino: 35 ejemplares para biblioteca y aulas, 15 para entrega simbólica (estudiantes destacados, aliados, autoridades), 10 para archivo institucional y muestras pedagógicas externas.

\subsection{Producto C · Revista \textbf{Cosmotec} (3 fascículos · financiamiento del semillero)}

\textbf{Insumo digital existente:} dos fascículos previos de la revista \textbf{Cosmotec} —ya diagramados, con contenidos producidos en ediciones anteriores del Centro de Interés— y un fascículo inédito en proceso de cierre editorial, articulado con la línea de astronomía (semillero Betelgeuse) y la línea de robótica e ingeniería (semillero CIR Sormatronik).

\textbf{Productos físicos a financiar:} \textbf{240 ejemplares totales (80 por fascículo $\times$ 3 fascículos)}, formato A4 (21{,}0\,\(\times\)\,29{,}7\,cm), entre 24 y 32 páginas, full color, encuadernación grapada al lomo, papel bond 90\,g (interior) y propalcote 200\,g (portada).

\textbf{Destino y modelo de circulación.} La revista \textbf{no se distribuye gratuitamente}: se entrega a la comunidad educativa, a familias, aliados y visitantes bajo un esquema de \textit{aporte voluntario / venta simbólica} —orden del 30\,\% al 50\,\% del costo unitario—. El producto íntegro de esa circulación se administra como \textbf{fondo del semillero institucional} y se destina a sufragar insumos de prototipos, transporte a ferias departamentales y nacionales, materiales de laboratorio y reposición de equipos básicos. De este modo, los \pesos{440.000} solicitados para Cosmotec funcionan como \textit{capital semilla editorial}: una inversión institucional que activa un mecanismo permanente de autosostenibilidad del semillero. Se reservan adicionalmente 12 ejemplares por fascículo para archivo institucional y muestras pedagógicas externas (no comercializables).

% ============================ PRESUPUESTO ============================
\section{Presupuesto solicitado — rubros tabulados}

\subsection{Tabla maestra de rubros}

\renewcommand{\arraystretch}{1.28}
\begin{table}[H]
\centering
\footnotesize
\rowcolors{2}{smjGrisClaro}{white}
\begin{tabularx}{\linewidth}{C{0.6cm} L{5.5cm} C{2.0cm} R{1.9cm} R{2.2cm}}
\toprule
\rowcolor{smjAzul}
\textcolor{white}{\textbf{N.°{}}} & \textcolor{white}{\textbf{Rubro}} & \textcolor{white}{\textbf{Cantidad}} & \textcolor{white}{\textbf{V. unit.}} & \textcolor{white}{\textbf{V. total (COP)}} \\
\midrule
1 & Impresión de \textbf{kits maestros de guías} (cartillas por grado-periodo de la plataforma ConectaTE; full color, anillado profesional, portada plastificada, $\sim$80 pág.\ c/u) & 24 cartillas\par(6 grados × 4 periodos) & \pesos{22.000} & \pesos{528.000} \\
2 & Producción de \textbf{bitácoras estudiantiles} (compilado de guías clave por grado, b/n, anillado sencillo, $\sim$20 pág.\ c/u) & 180 bitácoras\par(1 por estudiante) & \pesos{1.700} & \pesos{306.000} \\
3 & Edición e impresión del \textbf{libro «Salomé y los exoplanetas»} (A5, $\sim$60 pág., full color, encuadernación rústica cosida con solapas) — semillero Betelgeuse & 60 ejemplares & \pesos{9.000} & \pesos{540.000} \\
4 & Impresión de la \textbf{revista \emph{Cosmotec}} — 2 fascículos en reedición + 1 fascículo inédito (A4, 24--32 pág., full color, grapada al lomo, portada propalcote 200\,g) — \textit{financiamiento del semillero} & 240 ejemplares\par(80 × 3 fascículos) & \pesos{1.800} & \pesos{432.000} \\
5 & \textbf{Encuadernación profesional y acabados} (anillado doble-O, plastificado de portadas, rotulado institucional) & Servicio global & --- & \pesos{90.000} \\
6 & \textbf{Insumos editoriales} (cartulina Canson, separadores, fundas plásticas, papel propalcote de portadas) & Lote global & --- & \pesos{50.000} \\
7 & \textbf{Logística y distribución} (transporte desde la imprenta, embalaje, entrega en aulas y biblioteca) & Servicio global & --- & \pesos{30.000} \\
8 & \textbf{Imprevistos} (1{,}20\,\% del valor total — sobrecostos editoriales, reposiciones, variación de papel) & Reserva & --- & \pesos{24.000} \\
\midrule
\rowcolor{smjAzul!15}
\multicolumn{4}{r}{\textbf{\color{smjAzul} TOTAL SOLICITADO}} & \textbf{\color{smjAzul}\pesos{2.000.000}} \\
\bottomrule
\end{tabularx}
\caption{Rubros detallados del presupuesto solicitado para los tres productos editoriales del Centro de Interés ConectaTE.}
\end{table}

\begin{cajaTotal}
\centering
\textbf{\large\color{smjVino} Valor total solicitado:} \quad
{\Large\bfseries\color{smjAzul}\pesos{2.000.000 COP}}\\[2pt]
\small\itshape (dos millones de pesos m/cte — exactos)
\end{cajaTotal}

\subsection{Distribución porcentual por producto}

\renewcommand{\arraystretch}{1.25}
\begin{table}[H]
\centering
\small
\begin{tabularx}{\linewidth}{L{5.6cm} R{2.4cm} R{1.8cm} X}
\toprule
\rowcolor{smjVino!12}
\textbf{Producto / componente} & \textbf{Valor (COP)} & \textbf{\%} & \textbf{Justificación} \\
\midrule
A · Guías de la plataforma ConectaTE (rubros 1 y 2) & \pesos{834.000} & 41{,}70\,\% & Cubre los 6 grados y los 4 periodos del año lectivo \\
B · Libro Salomé — semillero Betelgeuse (rubro 3) & \pesos{540.000} & 27{,}00\,\% & Producto editorial STEAM+ del Centro de Interés \\
C · Revista Cosmotec — financiamiento del semillero (rubro 4) & \pesos{432.000} & 21{,}60\,\% & Capital semilla editorial: activa autosostenibilidad del semillero institucional \\
D · Acabados, insumos, logística e imprevistos (rubros 5--8) & \pesos{194.000} & 9{,}70\,\% & Soporte transversal a los productos A, B y C \\
\midrule
\rowcolor{smjGrisClaro}
\textbf{TOTAL} & \textbf{\pesos{2.000.000}} & \textbf{100\,\%} & \\
\bottomrule
\end{tabularx}
\caption{Distribución del presupuesto por producto editorial del Centro de Interés ConectaTE.}
\end{table}

\subsection{Modelo de retorno económico de la revista Cosmotec}

A diferencia de los productos A y B, la revista \textbf{Cosmotec} se concibe como \textit{producto editorial con retorno}. La proyección conservadora del primer año contempla:

\renewcommand{\arraystretch}{1.25}
\begin{table}[H]
\centering
\small
\begin{tabularx}{\linewidth}{L{5.6cm} R{2.4cm} R{2.4cm} X}
\toprule
\rowcolor{smjAzul!12}
\textbf{Escenario} & \textbf{Aporte simbólico} & \textbf{Ejemplares circulados} & \textbf{Recaudo proyectado (COP)} \\
\midrule
Escenario conservador (50\,\% del costo) & \pesos{900} c/u & 200 de 240 & \pesos{180.000} \\
Escenario base (aporte simbólico estándar) & \pesos{2.000} c/u & 200 de 240 & \pesos{400.000} \\
Escenario favorable (aporte sostén) & \pesos{3.000} c/u & 200 de 240 & \pesos{600.000} \\
\bottomrule
\end{tabularx}
\caption{Proyección de recaudo de la revista Cosmotec por escenario de aporte (240 ejemplares totales, $\sim$40 reservados para archivo).}
\end{table}

El recaudo se administra en una cuenta interna del semillero supervisada por el docente responsable y la rectoría, con rendición de cuentas pública al final del año lectivo.

\subsection{Cotizaciones de referencia}

Los valores unitarios fueron estimados con base en cotizaciones del sector gráfico de Cartago y norte del Valle (corredor Cartago–La\,Virginia–Pereira), conforme a tarifas vigentes en abril de 2026. El presupuesto contempla IVA incluido cuando aplica y reserva un margen del 2{,}25\,\% para variaciones del insumo de papel.

% ============================ CRONOGRAMA ============================
\section{Cronograma de ejecución}

\renewcommand{\arraystretch}{1.25}
\begin{table}[H]
\centering
\small
\begin{tabularx}{\linewidth}{C{1.4cm} L{5.2cm} C{2.0cm} X}
\toprule
\rowcolor{smjAzul!12}
\textbf{Semana} & \textbf{Actividad} & \textbf{Responsable} & \textbf{Producto} \\
\midrule
1 & Cierre de archivos digitales (PDFs finales de las 180 guías, maqueta del libro Salomé, fascículos 1 y 2 de Cosmotec listos, cierre del fascículo 3 inédito) y selección de imprenta. & Docente responsable + semillero & Archivos listos para imprenta \\
2 & Cotización formal, orden de servicio y pruebas de impresión (\emph{ozalid} y prueba de color). & Docente / Rectoría & Pruebas aprobadas \\
3 & Impresión, encuadernación y control de calidad. & Imprenta contratada & Lote terminado \\
4 & Recepción, rotulado institucional, entrega a la biblioteca y a las aulas; acto de lanzamiento del Centro de Interés ConectaTE. & Docente + Biblioteca + Semilleros & Distribución completada \\
\bottomrule
\end{tabularx}
\caption{Cronograma de ejecución (4 semanas desde la aprobación).}
\end{table}

% ============================ IMPACTO ============================
\section{Indicadores de impacto}

\renewcommand{\arraystretch}{1.30}
\begin{table}[H]
\centering
\small
\begin{tabularx}{\linewidth}{L{4.6cm} L{4.0cm} C{2.0cm} X}
\toprule
\rowcolor{smjVino!12}
\textbf{Indicador} & \textbf{Fórmula / fuente} & \textbf{Meta} & \textbf{Verificación} \\
\midrule
Cobertura estudiantil   & Estudiantes con bitácora / matriculados 6.°{}–11.°{} & $\geq$ 95\,\% & Lista de entrega por aula \\
Uso del kit maestro     & Sesiones documentadas con el kit en el cuaderno docente & $\geq$ 80\,\% de sesiones & Cuaderno docente \\
Lectura del libro Salomé & Estudiantes con bitácora de lectura diligenciada & $\geq$ 70\,\% & Bitácora institucional \\
Participación semilleros & Estudiantes activas/os en CIR Sormatronik o Betelgeuse & $\geq$ 30 estudiantes & Registro de semilleros \\
Circulación Cosmotec & Ejemplares con aporte recibido / 240 impresos & $\geq$ 80\,\% & Acta de cierre del fondo del semillero \\
Recaudo del semillero & Ingresos del fondo (escenario base) & $\geq$ \pesos{400.000} & Cuenta interna + informe a Rectoría \\
Producción estudiantil  & N.°{} de proyectos integradores entregados & $\geq$ 12 de 15 & Repositorio institucional \\
Satisfacción            & Encuesta Likert (1–5) a estudiantes y acudientes & $\geq$ 4{,}2 / 5 & Encuesta de cierre \\
\bottomrule
\end{tabularx}
\caption{Indicadores de impacto del Centro de Interés ConectaTE (año lectivo 2026).}
\end{table}

% ============================ SOSTENIBILIDAD ============================
\section{Sostenibilidad y proyección}

La inversión solicitada es \textbf{semilla, no recurrente}. El núcleo editorial digital del Centro de Interés \conectate{} (archivos YAML, pipeline de compilación, plataforma web, diagramación de Cosmotec) es propiedad intelectual del proyecto y queda disponible para reimpresiones futuras a costo marginal. Las guías impresas se diseñan con material durable (papel propalcote, plastificado, anillado doble-O) para soportar al menos \textbf{tres años lectivos consecutivos}. El libro «Salomé y los exoplanetas» se entrega como activo patrimonial de la biblioteca institucional.

Particularmente relevante es el caso de la revista \textbf{Cosmotec}: la inversión institucional opera como \textit{capital semilla editorial} que activa un mecanismo \textbf{autosostenible}. Los ingresos generados por la circulación de los tres fascículos en 2026 permitirán al semillero financiar autónomamente las ediciones de 2027 y 2028, sin requerir una nueva apropiación presupuestal de la institución. Es decir, los \pesos{432.000} aportados ahora se convierten —en proyección base— en un fondo recurrente del orden de \pesos{400.000} anuales para insumos, transporte y materiales del semillero.

Para escalar el Centro de Interés en 2027 y 2028, el docente responsable se compromete a:

\begin{itemize}[leftmargin=*]
  \item Presentar \conectate{} a convocatorias externas de fomento (Mincultura, Mineducación, Minciencias, Comfenalco, Computadores para Educar, fundaciones del Eje Cafetero).
  \item Mantener las guías digitales bajo control de versiones (\texttt{git}) y publicación abierta de la plataforma.
  \item Articular los semilleros CIR Sormatronik y Betelgeuse con redes externas (ferias de ciencia, encuentros departamentales, convocatorias nacionales).
  \item Documentar el proceso para que pueda ser replicado por otras áreas de la institución y por otras instituciones del municipio.
\end{itemize}

% ============================ COMPROMISOS ============================
\section{Compromisos del docente responsable}

\begin{enumerate}[label=\textbf{\color{smjAzul}C\arabic*.}, leftmargin=*]
  \item Custodiar y administrar los recursos asignados con criterios de transparencia, austeridad y trazabilidad.
  \item Presentar al Consejo Directivo un \textbf{informe ejecutivo de ejecución} dentro de los 30 días siguientes a la entrega del material, anexando facturas, fotografías y registro de distribución.
  \item Reservar al menos \textbf{10 ejemplares del libro Salomé}, un kit maestro completo de guías y \textbf{12 ejemplares de cada fascículo de Cosmotec} para el archivo institucional permanente de la I.E.\ Sor María Juliana.
  \item Reconocer expresamente a la institución y al financiador en cada producto editorial, en los actos públicos y en la plataforma digital del Centro de Interés.
  \item Donar los archivos fuente (\texttt{.yaml}, \texttt{.tex}, \texttt{.indd}/\texttt{.pdf} de Cosmotec) al repositorio institucional para garantizar la sostenibilidad y reimpresión futura.
  \item \textbf{Administrar con transparencia los recursos generados por la revista Cosmotec}, llevando una contabilidad sencilla del fondo del semillero (ingresos, egresos, saldo) y rindiendo cuentas al Consejo Directivo al cierre del año lectivo.
\end{enumerate}

% ============================ CIERRE ============================
\section{Cierre}

\begin{cajaCita}
\textit{«ConectaTE no es una plataforma con guías; es un Centro de Interés que produce plataformas, guías, libros, revistas, semilleros y comunidad. Su valor está en que transforma la curiosidad en investigación, la investigación en proyecto, el proyecto en experiencia formativa, la experiencia formativa en cultura institucional —y la cultura institucional, cuando madura, en autosostenibilidad. Pedir presupuesto para imprimir las guías de la plataforma, el libro Salomé y la revista Cosmotec es, ante todo, pedir presupuesto para que ese ecosistema cruce el umbral entre la pantalla y el papel, entre la curiosidad individual y la memoria escolar compartida, y entre la dependencia presupuestal y la economía propia del semillero.»}
\hfill\textemdash{} PhD. Álvaro Cárdenas Orozco
\end{cajaCita}

\vspace{1.4cm}

% ============================ FIRMAS ============================
\noindent
\begin{minipage}[t]{0.45\textwidth}\centering
\rule{\linewidth}{0.5pt}\\[3pt]
\textbf{\color{smjAzul} PhD. Álvaro Cárdenas Orozco}\\
\small Docente responsable\\
Centro de Interés ConectaTE\\
I.E. Sor María Juliana
\end{minipage}
\hfill
\begin{minipage}[t]{0.45\textwidth}\centering
\rule{\linewidth}{0.5pt}\\[3pt]
\textbf{\color{smjAzul} Aprobación de Rectoría}\\
\small Rector(a)\\
I.E. Sor María Juliana\\
Cartago — Valle del Cauca
\end{minipage}

\vspace{1.2cm}

\begin{center}
\small\color{gray}
Documento elaborado en Cartago — Valle del Cauca, \today.\\
Centro de Interés ConectaTE · I.E.\ Sor María Juliana
\end{center}

\end{document}
